\begin{abstract}
Public facial-symmetry tools take a single 2D photo, detect landmarks, reflect them
about a midline and always return a score, even when the photo cannot support one.
We build a system that accepts one face photo and returns a 2D geometric asymmetry
measure (global and per region) only when it is trustworthy, together with an
accept/caution/reject decision and its reasons. The method has three parts:
(i) asymmetry measured with fixed anatomical left/right pairs and a \emph{fitted}
symmetry axis, rather than nearest-neighbour matching about an assumed vertical midline;
(ii) a \emph{mirror test}, which runs the landmark detector on the photo and on its
mirror image and compares the two results, giving a label-free check of landmark
reliability and averaging out any left/right detector bias;
(iii) a quality gate on face size, sharpness, exposure, head pose and mirror
consistency, with thresholds set on a training split.
Since no dataset contains true symmetry labels, we evaluate against the same measure
computed from human-annotated landmarks (LFPW/300-W, 224 test images) and under 26
controlled perturbations (5{,}824 perturbed images).
The fitted-axis measure agrees far better with this annotation-based reference than a
re-implemented public baseline (Spearman $\rho$ \rhoMthree{} vs.\ \rhoMzero{}); the
baseline's KD-tree matching pairs \kdWrong\% of points with the wrong anatomical
partner. The mirror test ranks images by landmark error without ground truth
($\rho$ = \rhoMirrorNme{}), and mirror averaging cuts score drift under perturbation by
\driftReduction\%. The gate reduces silent failures (scores reported with error
$>$1\% of the inter-ocular distance) from \silentPoolOne\% to \silentPoolThree\%.
We also find that on in-the-wild photos the measured 2D asymmetry follows head yaw
almost one-to-one ($\rho$ = \rhoYawRef{}), even when computed from expert annotations.
This means pose gating is essential, not optional.
\end{abstract}
